% =============================================================
% Abstract (English translation of Abstrak), Lampiran 11.
% Title = translated thesis title, then the word ABSTRACT, then
% one paragraph, 200-300 words.
% =============================================================
\thispagestyle{empty}
\begin{center}
\textit{Thesis Title in English}\par
\vspace{6mm}
{\bfseries ABSTRACT\par}
\vspace{10mm}
\end{center}

\noindent The rapid growth of data volume across various sectors demands analytical approaches capable of automatically and accurately extracting patterns and knowledge. This research is motivated by that challenge, specifically within the domain of \textit{[application domain, e.g., anomaly detection/text classification/time-series forecasting]}, where conventional approaches have not been able to deliver optimal results in terms of either accuracy or computational efficiency. This study aims to design and evaluate an approach based on \textit{[method/algorithm name, e.g., deep learning, a Bayesian probabilistic model, or a metaheuristic algorithm]} that improves performance on the problem compared to a baseline method. The research methodology consists of data collection and preprocessing, model architecture design, training and hyperparameter tuning, and evaluation using a cross-validation scheme on the \textit{[dataset name/source]} dataset. Model performance is evaluated using \textit{[metric name, e.g., accuracy, F1-score, RMSE]} and compared against several state-of-the-art baseline methods. The results show that the proposed method achieves a performance improvement of \textit{[value/percentage]} over the baseline methods, with consistent results across several testing scenarios. These findings indicate that the proposed approach has the potential to be applied to similar problems at a larger scale, while also opening opportunities for further research on model architecture optimization and computational efficiency.

\vspace{6mm}
\noindent\textbf{Keywords:} keyword one, keyword two, keyword three, keyword four.
\clearpage
